\section{Introduction}

% Keep most of intro, just update these key sentences:

% ORIGINAL CONTRIBUTION PARAGRAPH - UPDATE TO:

This paper makes three primary contributions:

\begin{enumerate}
\item \textbf{Domain-adaptive pattern system}: We develop and validate a 17-pattern system that adapts to domain-specific linguistic characteristics, achieving 84.9\% F1-score across biology, history, mathematics, and literature domains—a 36.6 percentage point improvement over baseline methods.

\item \textbf{Cross-domain analysis}: We provide the first systematic analysis of domain-specific definition patterns in student notes, revealing three distinct styles: scientific nomenclature (biology), historical narrative (history), and academic formal (mathematics/literature). Our analysis demonstrates that domain-adaptive approaches significantly outperform domain-agnostic methods, reducing the domain performance gap by 82\%.

\item \textbf{Annotated dataset}: We release a dataset of 40 manually annotated student notes (10 per domain) containing 237 ground-truth definitions, enabling future research in educational NLP and student modeling. Statistical validation confirms all improvements are significant ($p < 0.05$) with large effect sizes (Cohen's $d > 2.0$).
\end{enumerate}

% ADD THIS PARAGRAPH after contributions:

Our evaluation demonstrates strong cross-domain generalization, with all four domains achieving F1-scores between 80.7\% and 90.3\%. Bootstrap confidence intervals confirm system reliability (overall 95\% CI: [82.3\%, 87.4\%]), making it suitable for deployment in real educational settings. The pattern-based approach proves both effective and interpretable, with clear mappings between linguistic patterns and domain-specific writing conventions.
